\documentclass[11pt,a4paper]{article}
\usepackage[utf8]{inputenc}
\usepackage[T1]{fontenc}
\usepackage[margin=0.72in]{geometry}
\usepackage{enumitem}
\usepackage{titlesec}
\usepackage{xcolor}
\usepackage{hyperref}

\definecolor{primarycolor}{RGB}{33,37,41}
\definecolor{secondarycolor}{RGB}{0,102,153}
\definecolor{accentcolor}{RGB}{90,99,108}
\hypersetup{colorlinks=true,urlcolor=secondarycolor,linkcolor=secondarycolor}
\titleformat{\section}{\large\bfseries\color{primarycolor}}{}{0em}{}
\titlespacing*{\section}{0pt}{1.1em}{0.4em}
\setlist[itemize]{leftmargin=1.3em,itemsep=0.15em,topsep=0.2em}
\pagestyle{empty}
\setlength{\parindent}{0pt}
\setlength{\parskip}{0.2em}

\begin{document}

\begin{center}
{\LARGE\bfseries\color{primarycolor} MARTA KOWAL-BRANDT}\\[0.25em]
{\large QA Automation Engineer | DevOps and Platform Engineering}\\[0.3em]
marta.kowalbrandt@example.com | +48 555 303 303 | gitlab.com/mkb-example
\end{center}

\section*{PROFESSIONAL SUMMARY}
QA automation engineer with 7 years of experience, including a recent focus on CI/CD pipelines, test environments and build tooling. Owns GitLab CI pipelines for three product teams and has built Helm-based ephemeral environments on a shared Kubernetes cluster. Personal experience includes a k3s homelab with GitOps, monitoring and Terraform-provisioned VPS infrastructure. Seeking to move into a DevOps or platform engineering role.

\section*{TECHNICAL SKILLS}
CI/CD: GitLab CI, Jenkins, pipeline optimization, artifact publishing\\
Containers and orchestration: Docker, Docker Compose, Helm, Kubernetes (namespace-level access and test-environment deployments), k3s\\
Infrastructure as Code and GitOps: Terraform (personal VPS project), Argo CD, Ansible (basic)\\
Observability: Prometheus, Grafana, Loki; test metrics dashboards\\
Cloud learning: AWS EC2 and S3 basics (personal learning)\\
Languages and data: Python, Bash, Java (test code), SQL, PostgreSQL\\
Testing: pytest, Playwright, Selenium, Postman\\
Systems: Linux, Git

\section*{SELECTED INFRASTRUCTURE PROJECTS}
\textbf{Merge-Request Ephemeral Environments | Odrava Health Software | 2023--2024}
\begin{itemize}
\item Built a GitLab CI workflow that deploys a short-lived application and PostgreSQL environment with seeded data to the shared Kubernetes cluster using Helm.
\item Configured automatic environment cleanup after merge or 48 hours of inactivity; QA and product teams use the environments for feature review and e2e testing.
\item Technologies: GitLab CI, Helm, Kubernetes, Docker, Bash, PostgreSQL.
\end{itemize}

\textbf{CI Pipeline Speed-up | Odrava Health Software | 2022}
\begin{itemize}
\item Profiled pipelines for three teams, parallelized test jobs, added dependency caching and separated e2e tests into their own stage.
\item Reduced average pipeline time from 38 to 14 minutes and monthly CI runner cost by about 30\%.
\item Technologies: GitLab CI, Docker, pytest-xdist.
\end{itemize}

\textbf{Flaky Test Dashboard | Odrava Health Software | 2023}
\begin{itemize}
\item Collected CI test results in PostgreSQL and built Grafana dashboards showing flaky tests by team.
\item Helped teams identify and address frequent offenders; the flaky test rate fell from 7\% to below 1\%.
\item Technologies: Python, PostgreSQL, Grafana, GitLab API.
\end{itemize}

\textbf{Homelab k3s Cluster with GitOps | Personal | 2024--2025}
\begin{itemize}
\item Set up a three-node Raspberry Pi 5 k3s cluster with applications deployed from a Git repository using Argo CD.
\item Added Prometheus, Grafana and Loki monitoring, Traefik ingress and cert-manager for TLS.
\item Provisioned infrastructure for a small Hetzner VPS with Terraform.
\item Technologies: k3s, Argo CD, Prometheus, Grafana, Loki, Terraform, Ansible (basic).
\end{itemize}

\section*{PROFESSIONAL EXPERIENCE}
\textbf{QA Automation Engineer (Test Infrastructure) | Odrava Health Software | 2021--Present}
\begin{itemize}
\item Own GitLab CI pipelines for three product teams (approximately 25 developers), covering builds, unit, integration and e2e test stages, and artifact publishing.
\item Built ephemeral test environments for merge requests using Docker Compose and later Helm charts on the company's shared Kubernetes cluster.
\item Maintain a pytest and Playwright e2e suite of approximately 900 tests; reduced the flaky rate from 7\% to below 1\%.
\item Wrote Bash and Python tooling for PostgreSQL test-data seeding.
\item Wrote the internal guide to how CI works, now used in team onboarding.
\item Mentored two QA interns in 2023.
\end{itemize}

\textbf{QA Engineer | Brightmill Software House | 2019--2021}
\begin{itemize}
\item Tested web and mobile applications through manual and automated testing for client projects.
\item Introduced a Selenium and Java automated regression suite for the largest client project.
\item Configured Jenkins jobs for nightly regression runs.
\end{itemize}

\textbf{Customer Support Specialist (2nd line) | Brightmill Software House | 2018--2019}
\begin{itemize}
\item Investigated customer issues using logs and databases, and reproduced bugs for developers.
\end{itemize}

\section*{EDUCATION AND CERTIFICATIONS}
BEng, Computer Engineering | Silesian Lowlands University of Technology | 2018\\
ISTQB Foundation Level | 2019\\
Certified Kubernetes Administrator (CKA) preparation; exam booked for January 2027 (not yet certified)

\section*{LANGUAGES}
Polish (native), English (C1), German (A2)

\end{document}